\begin{textsample}{NEG Dim 1 – global\_north\_2024\_12 – Score -1.00 – global\_north\_2024\_12/118888109.txt}  \label{ex:f1_neg_189}
Louise Adler on fear, freedom and the appalling state of public discourse Save articles for later We need to talk about the problems of race, identity and racism and stop being afraid to have complex, nuanced conversations, says Louise Adler, head of Adelaide Writers ' Week.
Louise Adler.Credit : Eddie Jim In her third event, to be held in March, Adler addresses these issues head-on, enlisting TV host and journalist Waleed Aly and academic Susan Carland, his wife, to talk about Islamophobia, as well as professor of history at Columbia Simon Schama on antisemitism.
Also on the bill is African-American writer Wesley Lowery on racism in America, Russian-American journalist Masha Gessen on Russia, Trump and Gaza, and Indian essayist Pankaj Mishra, who has written about the legacy of the Holocaust.
Writers ' Week is bound to be controversial.
Adler says she has brought together writers"who have thought more deeply than most of us can in our ordinary lives about issues, be they personal or public"and she hopes the audience goes away inspired."I would hope that as citizens of Australia and of the world we would care deeply about the horrors we are witnessing internationally, that our collective and individual voices would in fact would be loud and clear in defence of humanity,"Adler says."There has been a lot of vacuous rhetoric about the importance of social cohesion, that in this multicultural society, we should not import divisions from elsewhere.
Are we so self-absorbed, so complacent, so insular that world events -- whether in Ukraine, Syria, Gaza or the US -- should not be part of our public discourse?"Loading Held simultaneously with the Adelaide Festival, Writers ' Week has come under fire under Adler ' s stewardship for featuring Palestinian authors."The Israel lobby decided that it was reprehensible and that it didn't like those views and they didn't approve of them,"she told the ABC.
She argued that they have every right not to approve, but"they were not entitled to say those Palestinian writers should not therefore be included in a literary festival..."There were letters to the Adelaide Festival board, articles in the media, demands on sponsors to withdraw funding and pressure on the South Australian premier, she said :"Thankfully the Adelaide Festival board was resolute that this was an opportunity for people to hear Palestinians talk about their writing."Advertisement Adler described it is a tragedy that"the wish that we should have peace and justice and a self-determination for Palestinian people in the Middle East, that everybody should live in peace and security, that that ' s suppressed, that can't be discussed here".
The former publisher laments the state of public discourse in Australia, which she argues is demoralising and ill-informed :"The lack of analysis and seriousness, the hyperbolic language that ' s used in the public sphere, does not enhance the quality of the discussion or our understanding as citizens."Adelaide Writers ' Week is free and next year includes Helen Garner, Geraldine Brooks, Tim Winton and Anthony Horowitz.
The AWW debate with Annabel Crabb and David Marr will explore \textbf{Oscar} Wilde ' s idea that ' true friends stab you from the front '.

% matched lemmas: oscar
\end{textsample}
